\documentclass[11pt,letterpaper]{article}
\usepackage[utf8]{inputenc}
\usepackage[T1]{fontenc}
\usepackage{geometry}
\geometry{margin=2cm}
\usepackage{hyperref}
\usepackage{setspace}
\usepackage{siunitx}
\sisetup{separate-uncertainty=true, range-phrase=--}
\DeclareSIUnit\kcalmol{kcal\,mol^{-1}}
\DeclareSIUnit\angstrom{\text{\AA}}
\singlespacing

\begin{document}
\emergencystretch=2em

\begin{flushleft}
\textbf{Myke Vital Sao Temgoua} \\
Department of Physics, Faculty of Science \\
University of Yaound\'{e} I \\
P.O. Box 812, Yaound\'{e}, Cameroon \\
\href{mailto:myke-vital.sao@facsciences-uy1.cm}{myke-vital.sao@facsciences-uy1.cm} \\[1.5em]
September 11, 2026
\end{flushleft}

\begin{flushleft}
Prof.\ Kenneth M.\ Merz, Jr., Editor-in-Chief \\
\textit{Journal of Chemical Information and Modeling} \\
American Chemical Society
\end{flushleft}

\vspace{0.5em}

\noindent Dear Prof.\ Merz,

We submit our manuscript \textbf{``Estimand Divergence Between Static Docking and Short Molecular Dynamics as a Triage Filter for African Antimalarial Natural Products''} for consideration as an Article in the \textit{Journal of Chemical Information and Modeling} (JCIM).

Using a curated collection derived from 396 African natural products and 454 synthetic antimalarials, we evaluated a multi-target docking framework across \num{136} wild-type and mutant receptor states of PfDHFR and PfCRT. On a primary candidate cohort ($n=\num{12}$), a novel Resistance-Resilience Score (RRS) provided a five-tier stratification decoupling predicted wild-type binding strength from mutant-state score retention. Explicit-solvent molecular dynamics (MD) simulations revealed a directional divergence in \qty{87.5}{\percent} (7/8) of matched candidate--target states, wherein dynamic pocket relaxation compensated for static grid-docking penalties. This structural stress test demonstrates that static docking scores and short MD trajectories provide complementary, non-equivalent triage metrics for natural product lead optimization.

We present three major methodological contributions: (i) conceptualizing \textit{Estimand Divergence} between static docking RRS and explicit-solvent MD pose retention as a positive triage filter; (ii) establishing a denominator-unbiased RRS framework that excludes non-binding wild-type baselines ($|S_{\mathrm{WT}}| < \qty{5.0}{\kcalmol}$); and (iii) validating protocol robustness across an independent external sensitivity panel of 39 compounds (\num{312} docking complexes) confirmed by GNINA convolutional neural network (CNN) affinity scoring with \qty{100}{\percent} Class-A agreement. Grounded in the JCIM molecular dynamics reporting guidelines established by Soares et al. (2023), short explicit-solvent trajectories (10~ns) function effectively as secondary structural stress tests.

Code, source data, analysis scripts, and model inputs are available on GitHub under an open-source licence; a permanent Zenodo snapshot will be archived upon publication. Statistics are reported with effect sizes, permutation $p$-values, bootstrap confidence intervals, and Bonferroni-Holm adjustments.

\noindent\textbf{Suggested reviewers.}
\begin{enumerate}
  \item Dr.\ Fidele Ntie-Kang, University of Buea --- African natural-product computational discovery
  \item Dr.\ Kelly Chibale, University of Cape Town --- antimalarial drug discovery
  \item Dr.\ John D.\ Chodera, Memorial Sloan Kettering --- free-energy methods and MD best practices
\end{enumerate}

\noindent This manuscript is original, not under consideration elsewhere, and all co-authors have approved the submission.

\vspace{1em}

\noindent Sincerely,

\vspace{1em}

\noindent \textbf{Myke Vital Sao Temgoua} \\
(on behalf of all co-authors)

\end{document}
